\section{Der ungerade Anteil der e-Reihe: Leibniz und der Verzicht auf die geraden Glieder}
\label{sec:ug}

\subsection{Fragestellung, Leitthese und Status der Aussagen}
\label{sec:ug:frage}

Die Leibniz-Reihe $\sum_{n\ge 0}(-1)^n/(2n+1)$ läuft ausschließlich über \emph{ungerade} Indizes. Die Taylorreihe der e-Funktion zerfällt ebenfalls in einen Anteil mit geraden und einen Anteil mit ungeraden Potenzen. Dieses Kapitel klärt mit vollständigen Beweisen, in welchem präzisen Sinn der ungerade Anteil die e-Funktion allein festlegt, warum auf die geraden Glieder verzichtet werden kann, und in welchem Sinn der ungerade Anteil der \emph{bedeutsamere} ist. Drei Fragen leiten die Untersuchung:

\begin{enumerate}[label=\textbf{F\arabic*.}]
\item Bestimmt der ungerade Anteil der Reihe die Zahl $e$ (und $e^x$) allein?
\item Gilt die Umkehrung, bestimmen also die geraden Glieder allein $e$? Wenn nicht: woran liegt die Asymmetrie?
\item In welchen Bedeutungen von „Gewicht“ wiegt der ungerade Anteil schwerer, und in welchen nicht?
\end{enumerate}

\paragraph{Leitthese.}
\emph{Der ungerade Anteil ist der bedeutsamere Anteil der e-Reihe: Er ist die feinere Information (aus ihm folgt der gerade Anteil, aber nicht umgekehrt), er trägt allein die Orientierung (Vorzeichen, Drehsinn, Zeitrichtung), und er ist der Erzeuger, aus dem der gerade Anteil durch Differentiation entsteht. Leibniz hat mit den ungeraden Gliedern deshalb den strukturell gewichtigeren Teil gewählt.}

\paragraph{Wichtige Präzisierung des Wortes „schwerer“.}
Der ungerade Anteil ist \emph{nicht} der betragsmäßig größere. Für alle reellen $x$ gilt $|E_{\mathrm u}(x)|<E_{\mathrm g}(x)$ (Satz~\ref{satz:ug:gewicht}); auch in der Energie ist der gerade Anteil knapp größer. Die Behauptung „die ungeraden Glieder wiegen schwerer“ ist daher im Sinne von Größe oder Energie \emph{falsch} und wird hier nicht behauptet. Sie ist im Sinne von \emph{Informationsgehalt, Feinheit und Orientierung} wahr und wird in diesem Sinn bewiesen. Diese Unterscheidung macht die Leitthese erst mathematisch belastbar.

\paragraph{Status der Aussagen.}
Alle mit \emph{Satz}, \emph{Lemma} und \emph{Korollar} bezeichneten Aussagen werden vollständig bewiesen. \emph{Bemerkungen} ordnen ein. Historische Aussagen darüber, was Leibniz wusste oder beabsichtigte, sind keine mathematischen Sätze; sie stehen getrennt in Abschnitt~\ref{sec:ug:historie}. Numerische Werte wurden mit hoher Genauigkeit unabhängig nachgerechnet.

\subsection{Definitionen und Grundlagen}
\label{sec:ug:grund}

\begin{definition}[e-Reihe, gerader und ungerader Anteil]\label{def:ug:anteile}
Für $z\in\mathbb{C}$ sei
\begin{align}
E(z)&:=\sum_{n=0}^{\infty}\frac{z^n}{n!}, &
E_{\mathrm g}(z)&:=\sum_{k=0}^{\infty}\frac{z^{2k}}{(2k)!}, &
E_{\mathrm u}(z)&:=\sum_{k=0}^{\infty}\frac{z^{2k+1}}{(2k+1)!}.
\end{align}
Wir schreiben $e^z:=E(z)$ und $e:=E(1)$. $E_{\mathrm g}$ heißt \emph{gerader}, $E_{\mathrm u}$ \emph{ungerader Anteil}. Die Partialsummen des ungeraden Anteils sind $S_N:=\sum_{k=0}^{N}\frac{1}{(2k+1)!}$ (Wert bei $z=1$).
\end{definition}

\begin{lemma}[Grundeigenschaften]\label{lem:ug:grund}
\begin{enumerate}[label=(\alph*)]
\item Die Reihen $E,E_{\mathrm g},E_{\mathrm u}$ konvergieren für alle $z\in\mathbb{C}$ absolut.
\item $E(z+w)=E(z)E(w)$; insbesondere $E(z)E(-z)=1$.
\item $E=E_{\mathrm g}+E_{\mathrm u}$, $E_{\mathrm g}$ ist gerade, $E_{\mathrm u}$ ist ungerade, und
$E_{\mathrm g}(z)=\tfrac12\big(E(z)+E(-z)\big)$, $E_{\mathrm u}(z)=\tfrac12\big(E(z)-E(-z)\big)$.
Die Zerlegung einer Funktion in einen geraden und einen ungeraden Anteil ist eindeutig.
\item $E_{\mathrm g}(z)^2-E_{\mathrm u}(z)^2=1$.
\item Für reelles $x$ gilt $E(x)>0$, $E_{\mathrm g}(x)\ge 1$ (Gleichheit nur bei $x=0$), und $E_{\mathrm u}(x)$ hat das Vorzeichen von $x$. Ferner $E_{\mathrm g}(x)-E_{\mathrm u}(x)=E(-x)$.
\end{enumerate}
\end{lemma}

\begin{proof}
(a) Das Quotientenkriterium liefert $\frac{|z|^{n+1}/(n+1)!}{|z|^n/n!}=\frac{|z|}{n+1}\to0$; also konvergiert $\sum|z|^n/n!$. Teilreihen einer absolut konvergenten Reihe konvergieren absolut.

(b) Die Reihen $E(z)$ und $E(w)$ konvergieren absolut; nach dem Satz von Mertens ist ihr Cauchy-Produkt gleich dem Produkt der Summen \cite{rudin_analysis}. Der $n$-te Term des Cauchy-Produkts ist nach dem binomischen Lehrsatz
\[
\sum_{k=0}^{n}\frac{z^k}{k!}\frac{w^{n-k}}{(n-k)!}=\frac1{n!}\sum_{k=0}^n\binom nk z^kw^{n-k}=\frac{(z+w)^n}{n!}.
\]
Also $E(z)E(w)=E(z+w)$, und mit $w=-z$: $E(z)E(-z)=E(0)=1$.

(c) Wegen absoluter Konvergenz darf nach geraden und ungeraden Indizes umgeordnet werden: $E=E_{\mathrm g}+E_{\mathrm u}$. Da $(-z)^{2k}=z^{2k}$ und $(-z)^{2k+1}=-z^{2k+1}$, gilt $E(-z)=E_{\mathrm g}(z)-E_{\mathrm u}(z)$. Addition bzw.\ Subtraktion liefert die Formeln. Eindeutigkeit: Ist $f=g+u=g'+u'$ mit $g,g'$ gerade und $u,u'$ ungerade, so ist $g-g'=u'-u$ gerade und ungerade zugleich, also $0$.

(d) $E_{\mathrm g}^2-E_{\mathrm u}^2=(E_{\mathrm g}+E_{\mathrm u})(E_{\mathrm g}-E_{\mathrm u})=E(z)E(-z)=1$ nach (b), (c).

(e) $E(x)=E(x/2)^2\ge0$ nach (b), und $E(x)\ne0$ wegen $E(x)E(-x)=1$; also $E(x)>0$. $E_{\mathrm g}(x)$ ist eine Summe nichtnegativer Terme mit erstem Term $1$. Für $x>0$ sind alle Terme von $E_{\mathrm u}(x)$ positiv; $E_{\mathrm u}$ ist ungerade. Die letzte Gleichung ist (c).
\end{proof}

\begin{lemma}[Ableitung und Paritätswechsel]\label{lem:ug:abl}
Auf $\mathbb{R}$ gilt $E_{\mathrm u}'=E_{\mathrm g}$, $E_{\mathrm g}'=E_{\mathrm u}$ und $E'=E$. Folglich $E_{\mathrm u}''=E_{\mathrm u}$ und $E_{\mathrm g}''=E_{\mathrm g}$.
\end{lemma}

\begin{proof}
Potenzreihen mit Konvergenzradius $\infty$ sind im Inneren gliedweise differenzierbar \cite{rudin_analysis}. Es ist $\frac{d}{dx}\frac{x^{2k+1}}{(2k+1)!}=\frac{x^{2k}}{(2k)!}$, also $E_{\mathrm u}'=E_{\mathrm g}$. Ferner verschwindet die Ableitung des konstanten Gliedes $1$ von $E_{\mathrm g}$, und für $k\ge1$ ist $\frac{d}{dx}\frac{x^{2k}}{(2k)!}=\frac{x^{2k-1}}{(2k-1)!}$; mit $j=k-1$ folgt $E_{\mathrm g}'=\sum_{j\ge0}\frac{x^{2j+1}}{(2j+1)!}=E_{\mathrm u}$. Die Summe liefert $E'=E$.
\end{proof}

\begin{bemerkung}[Bilanz der „verschwindenden“ Glieder]\label{bem:ug:bilanz}
Lemma~\ref{lem:ug:abl} beantwortet die Ausgangsfrage der Arbeit (Beobachtung 49) in verfeinerter Form. Bei $E_{\mathrm g}\mapsto E_{\mathrm u}$ geht genau ein Glied verloren, nämlich das konstante Glied $1$. Bei $E_{\mathrm u}\mapsto E_{\mathrm g}$ geht \emph{kein} Glied verloren; im Gegenteil entsteht aus dem linearen Glied $x$ das neue konstante Glied $1$. Die Bilanz über beide Anteile ist ausgeglichen: Was der gerade Anteil verliert, gewinnt er aus dem ungeraden zurück. Die Parität wechselt bei jeder Ableitung, die Periode des Paritätswechsels ist $2$: $D^2E_{\mathrm u}=E_{\mathrm u}$, $D^2E_{\mathrm g}=E_{\mathrm g}$. Der ungerade Anteil ist damit für sich allein eine Eigenfunktion von $D^2$ zum Eigenwert $1$ und benötigt zur Reproduktion unter $D^2$ den geraden Anteil nicht.
\end{bemerkung}

\subsection{Hauptsatz I: Der ungerade Anteil legt e und den geraden Anteil fest}
\label{sec:ug:haupt1}

\begin{satz}[Rekonstruktion aus dem ungeraden Anteil]\label{satz:ug:rekon}
Für alle $x\in\mathbb{R}$ gilt
\begin{align}
E_{\mathrm g}(x)&=\sqrt{1+E_{\mathrm u}(x)^2}, \label{eq:ug:gerade}\\
e^{x}&=E_{\mathrm u}(x)+\sqrt{1+E_{\mathrm u}(x)^2}, \qquad e^{-x}=\sqrt{1+E_{\mathrm u}(x)^2}-E_{\mathrm u}(x). \label{eq:ug:exp}
\end{align}
Insbesondere gilt mit $s:=E_{\mathrm u}(1)=\sum_{k\ge0}\frac1{(2k+1)!}$:
\begin{equation}\label{eq:ug:e}
e=s+\sqrt{1+s^2},\qquad \frac1e=\sqrt{1+s^2}-s .
\end{equation}
\end{satz}

\begin{proof}
Nach Lemma~\ref{lem:ug:grund}(d) ist $E_{\mathrm g}^2=1+E_{\mathrm u}^2$, nach (e) ist $E_{\mathrm g}(x)\ge1>0$. Die positive Wurzel ist daher die einzige Möglichkeit, das gibt \eqref{eq:ug:gerade}. Mit $E=E_{\mathrm g}+E_{\mathrm u}$ und $E(-x)=E_{\mathrm g}-E_{\mathrm u}$ folgt \eqref{eq:ug:exp}, und $x=1$ liefert \eqref{eq:ug:e}.
\end{proof}

\begin{bemerkung}[Zahlenwert]
$s=1{,}17520119364380\ldots$, $\sqrt{1+s^2}=1{,}54308063481524\ldots$, Summe $=2{,}71828182845905\ldots=e$. Die geraden Glieder $\sum 1/(2k)!=1{,}5430806\ldots$ müssen zur Berechnung von $e$ nicht mehr aufsummiert werden: Sie sind eine algebraische Funktion der ungeraden Summe.
\end{bemerkung}

\begin{korollar}[Die geraden Koeffizienten folgen aus den ungeraden]\label{kor:ug:koeff}
Als Identität von Potenzreihen gilt
\[
\sum_{k\ge0}\frac{x^{2k}}{(2k)!}=\sqrt{1+\Big(\sum_{j\ge0}\frac{x^{2j+1}}{(2j+1)!}\Big)^{2}} ,
\]
d.\,h.\ jeder gerade Koeffizient $1/(2k)!$ ist ein Polynom in den ungeraden Koeffizienten $1,\tfrac1{3!},\tfrac1{5!},\dots$
\end{korollar}

\begin{proof}
Die rechte Seite ist die Komposition der um $0$ analytischen Funktion $u\mapsto\sqrt{1+u}=\sum_m\binom{1/2}{m}u^m$ mit der Potenzreihe $u=E_{\mathrm u}(x)^2$ (die bei $x=0$ verschwindet), also nach dem Kompositionssatz für Potenzreihen um $0$ in eine Potenzreihe entwickelbar. Nach Satz~\ref{satz:ug:rekon} stimmt sie mit $E_{\mathrm g}$ auf ganz $\mathbb{R}$ überein; die Koeffizienten einer Potenzreihe sind eindeutig bestimmt (Identitätssatz). Die Koeffizienten der Wurzelreihe sind Polynome in denen von $E_{\mathrm u}^2$, also in den ungeraden Koeffizienten.
\end{proof}

\begin{beispiel}[Nachrechnen der ersten geraden Koeffizienten]
Mit $E_{\mathrm u}=x+\frac{x^3}{6}+\frac{x^5}{120}+\dots$ ist $E_{\mathrm u}^2=x^2+\frac{x^4}{3}+\frac{2x^6}{45}+\dots$, $E_{\mathrm u}^4=x^4+\frac{2x^6}{3}+\dots$, $E_{\mathrm u}^6=x^6+\dots$. Aus $\sqrt{1+u}=1+\frac u2-\frac{u^2}{8}+\frac{u^3}{16}-\dots$ folgt
\begin{align*}
[x^2]&=\tfrac12=\tfrac1{2!},\\
[x^4]&=\tfrac12\cdot\tfrac13-\tfrac18=\tfrac16-\tfrac18=\tfrac1{24}=\tfrac1{4!},\\
[x^6]&=\tfrac12\cdot\tfrac2{45}-\tfrac18\cdot\tfrac23+\tfrac1{16}=\tfrac1{45}-\tfrac1{12}+\tfrac1{16}=\tfrac{16-60+45}{720}=\tfrac1{720}=\tfrac1{6!}.
\end{align*}
Die geraden Koeffizienten werden also tatsächlich aus den ungeraden erzeugt.
\end{beispiel}

\subsection{Hauptsatz II: Die Asymmetrie zwischen ungeradem und geradem Anteil}
\label{sec:ug:haupt2}

Die Umkehrung von Satz~\ref{satz:ug:rekon} gilt nicht. Der Grund ist eine reine Symmetrieeigenschaft, die zunächst ohne jede Reihenentwicklung formuliert wird.

\begin{lemma}[Spiegelungslemma]\label{lem:ug:spiegel}
Sei $D\subseteq\mathbb{R}$ symmetrisch ($D=-D$) und $f:D\to\mathbb{R}$ mit Zerlegung $f=f_{\mathrm g}+f_{\mathrm u}$. Dann gilt für alle $x\in D$:
\[
f(x)-f(-x)=2f_{\mathrm u}(x).
\]
Insbesondere: $f$ unterscheidet $x$ von $-x$ genau dann, wenn $f_{\mathrm u}(x)\neq0$. Ferner ist keine gerade Funktion auf einer symmetrischen Menge $D\ne\{0\}$ injektiv, während ungerade Funktionen injektiv sein können.
\end{lemma}

\begin{proof}
$f(x)-f(-x)=\big(f_{\mathrm g}(x)+f_{\mathrm u}(x)\big)-\big(f_{\mathrm g}(x)-f_{\mathrm u}(x)\big)=2f_{\mathrm u}(x)$. Für gerades $g$ und $x\ne0$ in $D$ ist $-x\in D$, $-x\ne x$ und $g(-x)=g(x)$, also ist $g$ nicht injektiv. Die Funktion $x\mapsto x$ ist ein Beispiel einer injektiven ungeraden Funktion.
\end{proof}

\begin{satz}[Verfeinerungssatz]\label{satz:ug:verfein}
Sei $\varphi(s):=\sqrt{1+s^2}$. Dann gilt:
\begin{enumerate}[label=(\alph*)]
\item $E_{\mathrm g}=\varphi\circ E_{\mathrm u}$ auf $\mathbb{R}$. Der gerade Anteil ist eine Funktion des ungeraden.
\item Es gibt \emph{keine} Funktion $\psi$ mit $E_{\mathrm u}=\psi\circ E_{\mathrm g}$ auf $\mathbb{R}$.
\item $E_{\mathrm u}:\mathbb{R}\to\mathbb{R}$ ist streng monoton wachsend und bijektiv, mit Umkehrfunktion $\operatorname{arsinh}(y)=\ln\!\big(y+\sqrt{1+y^2}\big)$.
\item $E_{\mathrm g}(x)=E_{\mathrm g}(-x)$; jeder Wert $c>1$ wird von $E_{\mathrm g}$ genau zweimal angenommen, nämlich in $\pm x$, und der Wert $1$ genau einmal (in $0$).
\end{enumerate}
Die Information, die $E_{\mathrm g}(x)$ über $x$ liefert, ist somit eine echte Vergröberung der Information von $E_{\mathrm u}(x)$.
\end{satz}

\begin{proof}
(a) ist \eqref{eq:ug:gerade}.

(b) Angenommen, $E_{\mathrm u}=\psi\circ E_{\mathrm g}$. Für $x\ne0$ ist $E_{\mathrm g}(x)=E_{\mathrm g}(-x)$, also $E_{\mathrm u}(x)=\psi(E_{\mathrm g}(x))=\psi(E_{\mathrm g}(-x))=E_{\mathrm u}(-x)=-E_{\mathrm u}(x)$, also $E_{\mathrm u}(x)=0$. Das widerspricht Lemma~\ref{lem:ug:grund}(e) ($E_{\mathrm u}(x)\ne0$ für $x\neq0$).

(c) Nach Lemma~\ref{lem:ug:abl} und~\ref{lem:ug:grund}(e) ist $E_{\mathrm u}'=E_{\mathrm g}\ge1>0$, also streng monoton. Für $x\ge0$ folgt $E_{\mathrm u}(x)=\int_0^xE_{\mathrm g}\ge x$; wegen der Ungeradheit ist $E_{\mathrm u}$ unbeschränkt nach oben und unten, und nach dem Zwischenwertsatz surjektiv. Ist $y=E_{\mathrm u}(x)$, so folgt aus \eqref{eq:ug:exp} $e^x=y+\sqrt{1+y^2}$, also $x=\ln(y+\sqrt{1+y^2})$.

(d) Die Gleichheit ist die Geradheit. Ferner ist $E_{\mathrm g}'=E_{\mathrm u}>0$ auf $(0,\infty)$, also ist $E_{\mathrm g}$ dort streng monoton wachsend mit $E_{\mathrm g}(0)=1$ und $E_{\mathrm g}(x)\ge1+\tfrac{x^2}{2}\to\infty$. Jeder Wert $c>1$ wird also auf $(0,\infty)$ genau einmal angenommen und auf $(-\infty,0)$ spiegelbildlich genau einmal.
\end{proof}

\begin{korollar}[Die geraden Glieder bestimmen $e$ nur bis auf Inversion]\label{kor:ug:invers}
Sei $c:=E_{\mathrm g}(1)=\sum_{k\ge0}\frac1{(2k)!}$. Dann gilt
\[
\big\{\,c+\sqrt{c^2-1},\;c-\sqrt{c^2-1}\,\big\}=\Big\{e,\tfrac1e\Big\},
\]
und aus $c$ allein lässt sich \emph{nicht} entscheiden, welche der beiden Zahlen $e=e^{+1}$ und $e^{-1}$ gemeint ist. Die Entscheidung erfordert genau ein zusätzliches Bit, das Vorzeichen des Arguments. Der ungerade Anteil liefert $e$ dagegen eindeutig (Gleichung~\eqref{eq:ug:e}).
\end{korollar}

\begin{proof}
Nach Lemma~\ref{lem:ug:grund}(d) ist $E_{\mathrm u}(1)^2=c^2-1$, also $E_{\mathrm u}(1)=\pm\sqrt{c^2-1}$ und, da $E_{\mathrm u}(1)>0$, gilt $E_{\mathrm u}(1)=+\sqrt{c^2-1}$ \emph{nur dann}, wenn man weiß, dass das Argument $+1$ und nicht $-1$ ist. Mit $E(1)=c+E_{\mathrm u}(1)$ und $E(-1)=c-E_{\mathrm u}(1)$ folgt die Mengengleichung. Da $E_{\mathrm g}(1)=E_{\mathrm g}(-1)$, ist $c$ für die Argumente $+1$ und $-1$ identisch.
\end{proof}

\begin{satz}[Informationsgehalt über die Orientierung]\label{satz:ug:info}
Sei $X$ eine reelle Zufallsvariable mit $P(X=0)=0$, deren Verteilung symmetrisch ist ($X$ und $-X$ haben dieselbe Verteilung). Dann sind $\operatorname{sgn}X$ und $E_{\mathrm g}(X)$ stochastisch unabhängig, und $\operatorname{sgn}X$ ist eine Funktion von $E_{\mathrm u}(X)$. Also
\[
H\big(\operatorname{sgn}X\mid E_{\mathrm g}(X)\big)=1\ \text{Bit},\qquad H\big(\operatorname{sgn}X\mid E_{\mathrm u}(X)\big)=0\ \text{Bit}.
\]
\end{satz}

\begin{proof}
Nach Satz~\ref{satz:ug:verfein}(d) ist $E_{\mathrm g}$ auf $[0,\infty)$ injektiv, also erzeugt $E_{\mathrm g}(X)=E_{\mathrm g}(|X|)$ dieselbe $\sigma$-Algebra wie $|X|$. Für jede Borelmenge $A\subseteq(0,\infty)$ gilt wegen der Symmetrie $P(X\in A)=P(X\in-A)$, also $P(X>0,\,|X|\in A)=\tfrac12P(|X|\in A)$. Damit ist $\operatorname{sgn}X$ unabhängig von $|X|$, mit $P(\operatorname{sgn}X=\pm1)=\tfrac12$. Daher $H(\operatorname{sgn}X\mid E_{\mathrm g}(X))=H(\operatorname{sgn}X)=1$ Bit. Nach Lemma~\ref{lem:ug:grund}(e) ist $\operatorname{sgn}E_{\mathrm u}(X)=\operatorname{sgn}X$, also ist die bedingte Entropie $0$.
\end{proof}

\begin{bemerkung}
Satz~\ref{satz:ug:info} macht die Aussage „der ungerade Anteil wiegt schwerer“ im Sinne des Informationsgehalts exakt: Der gerade Anteil enthält null Bit über die Orientierung, der ungerade Anteil enthält sie vollständig. Nach Lemma~\ref{lem:ug:spiegel} ist das kein Zufall der e-Funktion, sondern eine Folge der Symmetrie: \emph{Jede} Unterscheidung von $x$ und $-x$ steckt ausschließlich im ungeraden Anteil.
\end{bemerkung}

\subsection{Orientierung, Zeitrichtung und Stabilität im Tiefpass}
\label{sec:ug:zeit}

Für die Signaltheorie dieser Arbeit hat die Asymmetrie eine unmittelbare Bedeutung. Der Tiefpass hat die Lösung $e^{st}$ mit $s=-1/\tau$; das Vorzeichen von $s$ entscheidet zwischen Abklingen und Aufklingen.

\begin{satz}[Zeitrichtung und Stabilität stecken im ungeraden Anteil]\label{satz:ug:stab}
Für $s\in\mathbb{C}$ und $t\in\mathbb{R}$ gilt
\[
E_{\mathrm g}(st)=E_{\mathrm g}(-st),\qquad E_{\mathrm u}(st)=-E_{\mathrm u}(-st),\qquad e^{\pm st}=E_{\mathrm g}(st)\pm E_{\mathrm u}(st).
\]
Für $t>0$ und reelles $s\neq0$ haben also die Lösungen $e^{st}$ und $e^{-st}$ (eine ist abklingend, die andere anwachsend) \emph{denselben} geraden Anteil, und sie unterscheiden sich ausschließlich im Vorzeichen des ungeraden Anteils. Insbesondere gibt es keine Funktion des geraden Anteils $E_{\mathrm g}(st)$, die einen stabilen Pol $s<0$ von dem instabilen Pol $-s>0$ unterscheidet, während $E_{\mathrm u}(st)$ diese Unterscheidung trägt: $\operatorname{sgn}E_{\mathrm u}(st)=\operatorname{sgn}(s)$.
\end{satz}

\begin{proof}
Die ersten beiden Gleichungen sind Lemma~\ref{lem:ug:grund}(c), die dritte folgt aus $E(\pm z)=E_{\mathrm g}(z)\pm E_{\mathrm u}(z)$. Ein Widerspruch zur Existenz der Funktion: Wäre $F$ mit $F(E_{\mathrm g}(st))=1$ für $s<0$ und $=0$ für $-s>0$ definiert, folgte aus $E_{\mathrm g}(st)=E_{\mathrm g}((-s)t)$ sofort $1=0$. Die Vorzeichenaussage folgt aus Lemma~\ref{lem:ug:grund}(e) für $x=st$ mit $t>0$.
\end{proof}

\begin{korollar}[Abklingen und Aufklingen des RC-Glieds]\label{kor:ug:rc}
Für die Zeitkonstante $\tau>0$ und $t>0$ ist
\[
e^{-t/\tau}=E_{\mathrm g}(t/\tau)-E_{\mathrm u}(t/\tau),\qquad e^{+t/\tau}=E_{\mathrm g}(t/\tau)+E_{\mathrm u}(t/\tau).
\]
Das Abklingen ist die Differenz, das Aufklingen die Summe derselben beiden Anteile. Die Sprungantwort $1-e^{-t/\tau}=1-E_{\mathrm g}(t/\tau)+E_{\mathrm u}(t/\tau)$ ist monoton wachsend und beschränkt; im Vergleich zum Aufklingen kehrt allein das Vorzeichen des ungeraden Anteils die Richtung um.
\end{korollar}

\begin{proof}
Die Gleichungen folgen aus Satz~\ref{satz:ug:stab} mit $s=\mp1/\tau$. Zur letzten Aussage: $\frac{d}{dt}(1-e^{-t/\tau})=\frac1\tau e^{-t/\tau}>0$ und $0<e^{-t/\tau}<1$ für $t>0$.
\end{proof}

\subsection{Die Rolle des ungeraden Anteils in der Differentialgleichung}
\label{sec:ug:dgl}

\begin{satz}[Fundamentallösung]\label{satz:ug:fund}
Die Differentialgleichung $y''=y$ auf $\mathbb{R}$ besitzt zu den Anfangswerten $y(0)=0$, $y'(0)=1$ genau die Lösung $y=E_{\mathrm u}$. Jede Lösung hat die Gestalt
\[
y=y(0)\,E_{\mathrm g}+y'(0)\,E_{\mathrm u}=y(0)\,E_{\mathrm u}'+y'(0)\,E_{\mathrm u}.
\]
Der gesamte Lösungsraum wird also vom ungeraden Anteil und seiner Ableitung aufgespannt.
\end{satz}

\begin{proof}
Nach Lemma~\ref{lem:ug:abl} lösen $E_{\mathrm u}$ und $E_{\mathrm g}$ die Gleichung; $E_{\mathrm u}(0)=0$, $E_{\mathrm u}'(0)=E_{\mathrm g}(0)=1$, und $E_{\mathrm g}(0)=1$, $E_{\mathrm g}'(0)=E_{\mathrm u}(0)=0$. Sei $y$ irgendeine Lösung und $u:=y-y(0)E_{\mathrm g}-y'(0)E_{\mathrm u}$. Dann ist $u''=u$, $u(0)=u'(0)=0$. Setze $p:=u'+u$ und $q:=u'-u$. Es ist $p'=u''+u'=u+u'=p$ und $q'=u''-u'=u-u'=-q$. Daraus folgt $(pE(-t))'=p'E(-t)-pE(-t)=0$, also $pE(-t)\equiv p(0)=0$, und $(qE(t))'=q'E(t)+qE(t)=0$, also $qE(t)\equiv q(0)=0$. Da $E(\pm t)>0$, gilt $p=q=0$, also $u=(p-q)/2=0$. Die Eindeutigkeit der Anfangswertlösung folgt ebenso.
\end{proof}

\begin{bemerkung}
Der ungerade Anteil ist damit die \emph{Impulsantwort} des Systems $y''=y$ (Anfangswert Null, Anfangsgeschwindigkeit Eins). Aus ihr entsteht die zweite Grundlösung $E_{\mathrm g}$ durch Differentiation, in Übereinstimmung mit Satz~\ref{satz:ug:rekon}. Die Rolle „Erzeuger“ und „Erzeugtes“ ist somit nicht symmetrisch verteilt: $E_{\mathrm g}=E_{\mathrm u}'$, aber $E_{\mathrm u}$ ist nicht als Funktion von $E_{\mathrm g}$ darstellbar (Satz~\ref{satz:ug:verfein}\,(b)). Für die Auflösung der ersten Ordnung $y'=y$ werden beide Anteile gebraucht; in zweiter Ordnung genügt der ungerade Anteil samt Ableitung.
\end{bemerkung}

\subsection{Abgrenzung: Größe und Energie, in denen der gerade Anteil überwiegt}
\label{sec:ug:gewicht}

\begin{satz}[Gewichtsvergleich]\label{satz:ug:gewicht}
\begin{enumerate}[label=(\alph*)]
\item Für alle $x\in\mathbb{R}$ gilt $|E_{\mathrm u}(x)|<E_{\mathrm g}(x)$ und $E_{\mathrm g}(x)-|E_{\mathrm u}(x)|=e^{-|x|}$. Das Verhältnis $|E_{\mathrm u}|/E_{\mathrm g}$ wächst mit $|x|$ und strebt gegen $1$.
\item Auf $[-a,a]$, $a>0$, gilt mit dem $L^2$-Skalarprodukt $\langle f,g\rangle=\int_{-a}^{a}fg\,dx$:
\[
\|E_{\mathrm g}\|^2=a+\tfrac12E_{\mathrm u}(2a),\quad \|E_{\mathrm u}\|^2=-a+\tfrac12E_{\mathrm u}(2a),\quad\langle E_{\mathrm g},E_{\mathrm u}\rangle=0,\quad\|E\|^2=E_{\mathrm u}(2a).
\]
Der Energieanteil des ungeraden Anteils ist $\tfrac12-\dfrac{a}{E_{\mathrm u}(2a)}<\tfrac12$ und strebt für $a\to\infty$ gegen $\tfrac12$.
\end{enumerate}
\end{satz}

\begin{proof}
(a) Für $x\ge0$ gilt $E_{\mathrm g}(x)-E_{\mathrm u}(x)=E(-x)=e^{-x}>0$ nach Lemma~\ref{lem:ug:grund}(e); für $x<0$ ist $|E_{\mathrm u}(x)|=E_{\mathrm u}(-x)$ und $E_{\mathrm g}(x)=E_{\mathrm g}(-x)$. Das Verhältnis ist $\tanh|x|=\frac{1-e^{-2|x|}}{1+e^{-2|x|}}$, denn $\frac{E_{\mathrm u}}{E_{\mathrm g}}=\frac{E(x)-E(-x)}{E(x)+E(-x)}=\frac{1-E(-2x)}{1+E(-2x)}$; das ist für $x>0$ streng wachsend in $x$ und konvergiert gegen $1$.

(b) Aus Lemma~\ref{lem:ug:grund}(c) folgt $E_{\mathrm g}^2=\tfrac14\big(E(2x)+2+E(-2x)\big)=\tfrac12\big(E_{\mathrm g}(2x)+1\big)$ und analog $E_{\mathrm u}^2=\tfrac12\big(E_{\mathrm g}(2x)-1\big)$. Wegen $\frac{d}{dx}E_{\mathrm u}(2x)=2E_{\mathrm g}(2x)$ ist $\int_{-a}^{a}E_{\mathrm g}(2x)\,dx=E_{\mathrm u}(2a)$. Also $\|E_{\mathrm g}\|^2=\tfrac12(E_{\mathrm u}(2a)+2a)$ und $\|E_{\mathrm u}\|^2=\tfrac12(E_{\mathrm u}(2a)-2a)$. Das Integral des ungeraden Produkts $E_{\mathrm g}E_{\mathrm u}$ über das symmetrische Intervall ist $0$. Zusammen: $\|E\|^2=E_{\mathrm u}(2a)$ (Kontrolle: $\int_{-a}^ae^{2x}dx=\tfrac12(e^{2a}-e^{-2a})=E_{\mathrm u}(2a)$). Der Anteil ist $\frac{\frac12E_{\mathrm u}(2a)-a}{E_{\mathrm u}(2a)}=\tfrac12-\frac a{E_{\mathrm u}(2a)}$; wegen $E_{\mathrm u}(2a)>2a$ ist er positiv und kleiner als $\tfrac12$, und $E_{\mathrm u}(2a)\ge\frac{(2a)^3}{6}$ liefert $a/E_{\mathrm u}(2a)\to0$.
\end{proof}

\begin{bemerkung}
Damit ist die Aussage „der ungerade Anteil ist der betragsmäßig oder energetisch schwerere“ \emph{widerlegt}; der gerade Anteil überwiegt in beiden Maßen, im Betrag um genau $e^{-|x|}$, in der Energie um genau $a/E_{\mathrm u}(2a)$ des Gesamtwerts. Beide Differenzen verschwinden asymptotisch. Das Gewicht des ungeraden Anteils liegt somit nicht in seiner Größe, sondern in der Information, die er trägt (Abschnitte~\ref{sec:ug:haupt1} und \ref{sec:ug:haupt2}). Die Unterscheidung ist wesentlich: Wer „schwerer“ als „größer“ liest, behauptet etwas Falsches; wer „schwerer“ als „informationsreicher, feiner und erzeugend“ liest, behauptet etwas Bewiesenes.
\end{bemerkung}

\subsection{Rechenökonomie: Wie viele Glieder braucht die Rekonstruktion?}
\label{sec:ug:rechnen}

\begin{satz}[Fehlerabschätzung der ungeraden Rekonstruktion]\label{satz:ug:fehler}
Sei $g(s):=s+\sqrt{1+s^2}$ und $r_N:=E_{\mathrm u}(1)-S_N=\sum_{k>N}\frac1{(2k+1)!}$. Dann gilt für alle $N\ge0$
\begin{align*}
0<r_N&\le\frac{16}{15}\cdot\frac1{(2N+3)!},\\
0<e-g(S_N)&<(1+\tanh 1)\,r_N\le\frac{16\,(1+\tanh1)}{15\,(2N+3)!}<\frac{1{,}88}{(2N+3)!}.
\end{align*}
\end{satz}

\begin{proof}
Es ist $(2N+3+2j)!=(2N+3)!\prod_{m=1}^{2j}(2N+3+m)\ge(2N+3)!\,(2N+4)^{2j}$, also
\begin{align*}
r_N&=\sum_{j\ge0}\frac1{(2N+3+2j)!}\le\frac1{(2N+3)!}\sum_{j\ge0}(2N+4)^{-2j}\\
&=\frac{1}{(2N+3)!}\cdot\frac1{1-(2N+4)^{-2}}\le\frac{16}{15}\cdot\frac1{(2N+3)!}.
\end{align*}
Die Funktion $g$ ist glatt mit $g'(s)=1+\frac s{\sqrt{1+s^2}}$ und $g''(s)=(1+s^2)^{-3/2}>0$, also ist $g'$ wachsend. Nach dem Mittelwertsatz gilt $g(s)-g(S_N)=g'(\xi)\,r_N$ mit $\xi\in(S_N,s)$, wobei $s=E_{\mathrm u}(1)$, und somit $g'(\xi)<g'(s)=1+\frac{E_{\mathrm u}(1)}{E_{\mathrm g}(1)}=1+\tanh1<1{,}7616$. Wegen $g(s)=e$ (Satz~\ref{satz:ug:rekon}) folgt die Behauptung; die letzte Abschätzung ist $\tfrac{16}{15}\cdot1{,}7616<1{,}88$.
\end{proof}

\begin{beobachtung}[Numerische Überprüfung]\label{beob:ug:tab}
Die folgende Tabelle wurde mit 40-stelliger Arithmetik berechnet. Spalte $e-g(S_N)$ ist der tatsächliche Fehler der ungeraden Rekonstruktion (jeweils $N+1$ Glieder), „Schranke“ die Abschätzung aus Satz~\ref{satz:ug:fehler}, „gerade Variante“ der Fehler von $C_N+\sqrt{C_N^2-1}$ mit $C_N=\sum_{k=0}^N\frac1{(2k)!}$ (ebenfalls $N+1$ Glieder, mit als bekannt vorausgesetztem Vorzeichen des Arguments).
\begin{center}
\small
\begin{tabular}{c c c c c c}
$N$ & $S_N$ & $g(S_N)$ & $e-g(S_N)$ & Schranke & gerade Variante\\\hline
0 & 1{,}00000000000 & 2{,}41421356237 & $3{,}04\cdot10^{-1}$ & $3{,}13\cdot10^{-1}$ & $1{,}72$\\
1 & 1{,}16666666667 & 2{,}70325740955 & $1{,}50\cdot10^{-2}$ & $1{,}57\cdot10^{-2}$ & $1{,}00\cdot10^{-1}$\\
2 & 1{,}17500000000 & 2{,}71792741242 & $3{,}54\cdot10^{-4}$ & $3{,}73\cdot10^{-4}$ & $3{,}27\cdot10^{-3}$\\
3 & 1{,}17519841270 & 2{,}71827692956 & $4{,}90\cdot10^{-6}$ & $5{,}18\cdot10^{-6}$ & $5{,}80\cdot10^{-5}$\\
4 & 1{,}17520116843 & 2{,}71828178404 & $4{,}44\cdot10^{-8}$ & $4{,}71\cdot10^{-8}$ & $6{,}42\cdot10^{-7}$\\
5 & 1{,}17520119348 & 2{,}71828182817 & $2{,}84\cdot10^{-10}$ & $3{,}02\cdot10^{-10}$ & $4{,}86\cdot10^{-9}$\\
\end{tabular}
\end{center}
Die Schranke wird stets eingehalten und ist bis auf etwa $6\,\%$ scharf. Bei gleicher Gliederzahl ist die ungerade Rekonstruktion um etwa eine Größenordnung genauer als die gerade: Die ungeraden Glieder $1/(2k+1)!$ fallen schneller als die geraden $1/(2k)!$, und die Wurzelfunktion der geraden Variante ist bei $c\to1$ schlecht konditioniert. Diese Beobachtung ist eine \emph{Rechenökonomie}, nicht der Kern der Asymmetrie; der Kern ist Satz~\ref{satz:ug:verfein}.
\end{beobachtung}

\subsection{Leibniz: Phase und Drehsinn im ungeraden Anteil}
\label{sec:ug:leibniz}

Die Leibniz-Reihe ist nicht der ungerade Anteil der e-Reihe selbst (deren Nenner sind Fakultäten), sondern der ungerade Anteil der Reihe der \emph{Umkehrfunktion}, des Logarithmus, an der Stelle $z=i$. Diese Beziehung ist streng beweisbar und zeigt, dass dieselbe Asymmetrie wie in Abschnitt~\ref{sec:ug:haupt2} dort wieder auftritt.

\begin{satz}[Leibniz-Reihe in der Logarithmusreihe]\label{satz:ug:log}
Sei $\ell(z):=\sum_{n\ge1}\frac{(-1)^{n+1}}{n}z^n$ (Radius $1$). Dann konvergiert $\ell(i)$, und
\[
\ell(i)=\underbrace{\sum_{k\ge1}\frac{(-1)^{k+1}}{2k}}_{\text{gerade Indizes}}+i\underbrace{\sum_{k\ge0}\frac{(-1)^k}{2k+1}}_{\text{ungerade Indizes: Leibniz}}=\tfrac12\ln2+i\,\frac\pi4=\operatorname{Log}(1+i).
\]
Insbesondere ist $\sum_{k\ge1}\frac{(-1)^{k+1}}{2k}=\ln\sqrt2=\ln|1+i|$ (Betrag) und die Leibniz-Reihe $L=\frac\pi4=\operatorname{Arg}(1+i)$ (Phase).
\end{satz}

\begin{proof}
Für gerades $n=2k$ ($k\ge1$) ist $\frac{(-1)^{2k+1}i^{2k}}{2k}=\frac{-(-1)^k}{2k}=\frac{(-1)^{k+1}}{2k}$ (reell). Für ungerades $n=2k+1$ ($k\ge0$) ist $\frac{(-1)^{2k+2}i^{2k+1}}{2k+1}=i\frac{(-1)^k}{2k+1}$ (rein imaginär). Realteil und Imaginärteil sind alternierende Reihen mit monoton fallenden Beträgen gegen $0$ und konvergieren nach dem Leibniz-Kriterium; also konvergiert $\ell(i)$. Für $|z|<1$ ist $\ell(z)=\operatorname{Log}(1+z)$ (Hauptzweig). Der Abelsche Grenzwertsatz für Potenzreihen \cite{apostol_calculus} liefert $\ell(i)=\lim_{r\to1^-}\ell(ri)=\lim_{r\to1^-}\operatorname{Log}(1+ri)=\operatorname{Log}(1+i)$, denn $\operatorname{Log}$ ist auf $\{\operatorname{Re}w>0\}\ni1+i$ stetig. Mit $\operatorname{Log}(1+i)=\ln\sqrt2+i\,\frac\pi4$ folgt der Vergleich von Real- und Imaginärteil.
\end{proof}

\begin{satz}[Spiegelungsverhalten]\label{satz:ug:konj}
Es gilt $\ell(\bar z)=\overline{\ell(z)}$ für $|z|\le1$, $z\ne-1$ (bei Konvergenz). Daher ist $\ell(-i)=\tfrac12\ln2-i\,\frac\pi4$: Der Anteil über gerade Indizes ($\ln\sqrt2$) ist für $1+i$ und $1-i$ \emph{identisch}, der Anteil über ungerade Indizes (Leibniz) wechselt das Vorzeichen. Die Leibniz-Reihe ist also der Anteil, der die beiden gespiegelten Punkte $1\pm i$ (Drehsinn $+45^\circ$ bzw.\ $-45^\circ$) voneinander unterscheidet; der gerade Anteil sieht nur den Betrag $\sqrt2$.
\end{satz}

\begin{proof}
Die Koeffizienten von $\ell$ sind reell, also $\ell(\bar z)=\sum\frac{(-1)^{n+1}}{n}\bar z^n=\overline{\ell(z)}$. Für $z=i$ folgt $\ell(-i)=\overline{\ell(i)}=\tfrac12\ln2-i\frac\pi4$ mit Satz~\ref{satz:ug:log}.
\end{proof}

\begin{korollar}[Betrag und Phase in der e-Sprache]\label{kor:ug:polar}
Mit $L:=\sum_{k\ge0}\frac{(-1)^k}{2k+1}=\frac\pi4$ gilt
\[
E\!\Big(\tfrac12\ln2\Big)=\sqrt2,\qquad E(iL)=\frac{1+i}{\sqrt2},\qquad E\big(\ell(i)\big)=1+i,
\]
und $E(4iL)=-1$. Der gerade Anteil der Logarithmusreihe liefert den Betrag $\sqrt2$, der Leibniz-Anteil liefert die Phase.
\end{korollar}

\begin{proof}
$E(a+ib)=E(a)E(ib)$ mit $E(ib)=\cos b+i\sin b$ (Lemma~\ref{lem:ug:grund}(c) mit $z=ib$: $E_{\mathrm g}(ib)=\sum(-1)^kb^{2k}/(2k)!=\cos b$ und $E_{\mathrm u}(ib)=i\sin b$). Mit $a=\ln\sqrt2$ und $b=\frac\pi4$ folgt $E(\ell(i))=\sqrt2\big(\tfrac1{\sqrt2}+\tfrac i{\sqrt2}\big)=1+i$. Ferner $E(4iL)=E(i\pi)=\cos\pi+i\sin\pi=-1$.
\end{proof}

\begin{satz}[Der Leibniz-Wert als Gleichgewicht von geradem und ungeradem Anteil]\label{satz:ug:gleichg}
Auf der imaginären Achse gilt $E_{\mathrm g}(i\theta)=\cos\theta$ und $E_{\mathrm u}(i\theta)=i\sin\theta$. Der kleinste positive Winkel, bei dem beide Anteile betragsgleich sind, $|E_{\mathrm g}(i\theta)|=|E_{\mathrm u}(i\theta)|$, ist $\theta=\frac\pi4=L$, der Wert der Leibniz-Reihe. Dort gilt $E_{\mathrm g}(iL)=\frac1{\sqrt2}$ und $E_{\mathrm u}(iL)=\frac{i}{\sqrt2}$.
\end{satz}

\begin{proof}
Die Darstellung folgt aus Lemma~\ref{lem:ug:grund} wie im Beweis von Korollar~\ref{kor:ug:polar}. Aus $|\cos\theta|=|\sin\theta|$ folgt $\cos^2\theta=\sin^2\theta$, d.\,h.\ $\cos2\theta=0$; die kleinste positive Lösung ist $2\theta=\frac\pi2$. Nach Satz~\ref{satz:ug:log} (oder dem Konvergenzsatz der Leibniz-Reihe weiter oben in dieser Arbeit) ist $L=\frac\pi4$.
\end{proof}

\begin{bemerkung}[Reelle und imaginäre Achse im Vergleich]
Auf der reellen Achse überwiegt der gerade Anteil stets um $e^{-|x|}$ (Satz~\ref{satz:ug:gewicht}); die Anteile gleichen sich nur asymptotisch. Auf der imaginären Achse wechseln die Anteile dagegen periodisch, und der Leibniz-Wert $\frac\pi4$ ist der erste Punkt des exakten Gleichgewichts. Hier liegt der mathematisch präzise Kern der Aussage in der Arbeit, dass sich bei $\pi/4$ Sinus und Cosinus, „Ungerades und Gerades“, treffen.
\end{bemerkung}

\begin{bemerkung}[Halbkreis]
Auch auf dem Kreis zeigt sich die Asymmetrie von Lemma~\ref{lem:ug:spiegel}: Auf dem symmetrischen Bereich $\theta\in(-\frac\pi2,\frac\pi2)$ (rechter Halbkreis) ist $\sin$ bijektiv auf $(-1,1)$, $\cos$ aber gerade und daher nicht injektiv. Dort gilt $E(i\theta)=\sqrt{1-\sin^2\theta}+i\sin\theta$, denn $\cos\theta>0$ auf diesem Bereich; der ungerade Anteil rekonstruiert den Punkt auf dem Kreis vollständig, der gerade nur bis auf Spiegelung. Der Wert $\frac\pi4$ liegt in der Mitte dieses Halbkreises.
\end{bemerkung}

\subsection{Historische und interpretative Einordnung}
\label{sec:ug:historie}

Die folgenden Punkte sind \emph{keine} mathematischen Sätze und werden nicht bewiesen; sie sind von den Ergebnissen dieses Kapitels logisch unabhängig.

\begin{enumerate}
\item Die Reihe $\sum(-1)^n/(2n+1)=\pi/4$ wird in der Literatur Leibniz zugeschrieben (17.~Jahrhundert); Vorläufer sind Gregory und die indische Madhava-Schule. Die Zuschreibung ist hier ohne Belang: Bewiesen ist die Aussage über die \emph{Reihe}, nicht über ihre Urheberschaft.
\item Die Behauptung der Arbeit, Leibniz habe die Zerlegung der e-Reihe nach geraden und ungeraden Gliedern gekannt und bewusst die ungeraden gewählt, ist eine historische These. Sie müsste durch Quellen (Briefe, Manuskripte) belegt werden und ist mit mathematischen Mitteln weder beweisbar noch widerlegbar.
\item Was sich mathematisch sagen lässt, ist: \emph{Wer die Struktur einer symmetrischen Funktion oder Reihe erfassen will, muss den ungeraden Anteil betrachten}, denn nur er unterscheidet $x$ von $-x$ (Lemma~\ref{lem:ug:spiegel}). Leibniz' Reihe ist genau dieser Anteil der Logarithmusreihe bei $z=i$ (Satz~\ref{satz:ug:log}). In diesem streng definierten Sinn hat die Leibniz-Reihe den bedeutsameren Teil gewählt, gleichgültig, ob dies Absicht war.
\end{enumerate}

\subsection{Ergebnis}
\label{sec:ug:ergebnis}

\begin{center}
\small
\begin{tabular}{p{0.34\textwidth} p{0.2\textwidth} p{0.34\textwidth}}
\textbf{Bedeutung von „Gewicht“} & \textbf{Überwiegt} & \textbf{Beleg}\\\hline
Betrag $|f(x)|$ & gerade (um $e^{-|x|}$) & Satz~\ref{satz:ug:gewicht}(a)\\
Energie ($L^2$) & gerade (knapp) & Satz~\ref{satz:ug:gewicht}(b)\\
Feinheit der Information & \textbf{ungerade} & Satz~\ref{satz:ug:verfein}, Kor.~\ref{kor:ug:invers}\\
Orientierung, Vorzeichen, Zeitrichtung, Stabilität & \textbf{ungerade, ausschließlich} & Lemma~\ref{lem:ug:spiegel}, Satz~\ref{satz:ug:info}, \ref{satz:ug:stab}\\
Erzeugung (Fundamentallösung, Ableitung) & \textbf{ungerade} & Satz~\ref{satz:ug:fund}\\
Umkehrbarkeit ($\operatorname{arsinh}$) & \textbf{ungerade} & Satz~\ref{satz:ug:verfein}(c)\\
Rechenökonomie der Rekonstruktion von $e$ & \textbf{ungerade} & Satz~\ref{satz:ug:fehler}, Beobachtung~\ref{beob:ug:tab}\\
Phase statt Betrag (Logarithmusreihe bei $i$) & \textbf{ungerade (Leibniz)} & Satz~\ref{satz:ug:log}, \ref{satz:ug:konj}
\end{tabular}
\end{center}

\paragraph{Zusammenfassung.}
Die Zahl $e$ ist durch den ungeraden Anteil der Taylorreihe vollständig und eindeutig festgelegt,
\[
e=\sum_{k\ge0}\frac1{(2k+1)!}+\sqrt{1+\Big(\sum_{k\ge0}\frac1{(2k+1)!}\Big)^{2}},
\]
und die geraden Glieder sind eine algebraische Folge davon; auf ihre Summation kann verzichtet werden. Die Umkehrung gilt nicht: Der gerade Anteil bestimmt $e$ nur bis auf den Kehrwert $1/e$ und benötigt das Vorzeichen des Arguments, also eine Information, die ausschließlich im ungeraden Anteil liegt. Der ungerade Anteil ist der \emph{bedeutsamere}: Er ist die feinere Information, der Träger von Orientierung und Zeitrichtung, der Erzeuger der zweiten Grundlösung und derjenige Anteil, in dem Leibniz' Reihe die Phase $\pi/4$ der Logarithmusreihe liefert. Dass er dabei nicht der betragsmäßig größere ist, sondern um $e^{-|x|}$ hinter dem geraden zurückbleibt, schwächt diese Aussage nicht, sondern präzisiert sie: Sein Gewicht ist das Gewicht der Information, nicht der Größe.

